% SHARD-ID: AQM-20-GRASSMANN-WEYL-FERMION-DISPLACEMENTS
% SHARD-TITLE: Grassmann-Weyl Fermion Displacements
% SHARD-SUMMARY: Defines the abstract Grassmann displacement ray algebra used for fermion coherent states.
% SHARD-SUMMARY: States the Cahill--Glauber representation inside the Grassmann-extended finite CAR algebra.
% SHARD-SUMMARY: Sets up the Lamport dependencies proved in the following derivation shard.
% SHARD-KEYWORDS: Grassmann, Weyl, CAR, fermions, coherent states, displacement operators, projective representation

\section{Grassmann-Weyl Fermion Displacements}
\label{sec:grassmann-weyl-fermion-displacements}

\subsection{Status and Source Anchors}

This shard corrects the Grassmann question left open by
Section~\ref{sec:canonical-boson-fermion-field-comparison}.  The intended
object is not a homomorphism sending a Grassmann generator to a CAR
annihilator.  The intended object is the fermionic displacement operator of
Cahill--Glauber:
\[
  D(\gamma)=
  \exp\!\left(\sum_i(a_i^\dagger\gamma_i-\gamma_i^*a_i)\right),
\]
where the \(a_i,a_i^\dagger\) are concrete CAR operators and the
\(\gamma_i,\gamma_i^*\) are Grassmann coefficients.

The canonical local source is
\path{references/canonical_fields/PhysRevA.59.1538.pdf}, converted with
\texttt{marker} to
\path{references/canonical_fields/physrevA_59_1538_marker/PhysRevA.59.1538/PhysRevA.59.1538.md}.
The marker line anchors are: lines 37--45 for the fermion CAR and vacuum,
49--66 for Grassmann variables, cross anticommutation, and the adjoint
convention; 72--81 for the definition and finite expansion of \(D(\gamma)\);
83--101 for the displacement of \(a_n\) and \(a_n^\dagger\); 103--133 for
normal ordering and the ray-composition law; 350--362 for displacement
operator completeness; and 364--368 for the caveat that physical quantum
operators themselves do not involve Grassmann variables.

\subsection{Lamport Dependency Skeleton}

\[
\begin{array}{rcl}
D1&:&I\text{ is a finite mode set.}\\
D2&:&\operatorname{CAR}(I)\text{ is generated by odd }a_i,a_i^\dagger
      \text{ with }\{a_i,a_j^\dagger\}=\delta_{ij}.\\
D3&:&\Lambda_I\text{ is generated by odd }\gamma_i,\gamma_i^*
      \text{ with all Grassmann anticommutators zero.}\\
D4&:&\mathcal A_I=\operatorname{CAR}(I)\widehat\otimes\Lambda_I
      \text{ is the graded tensor product.}\\
D5&:&\sigma(\alpha,\beta)=\sum_i(\beta_i^*\alpha_i-\alpha_i^*\beta_i).\\
L1&:&\sigma\text{ is an even central skew-adjoint Grassmann cocycle form.}\\
D6&:&\mathsf D_I\text{ is the abstract Grassmann displacement ray algebra.}\\
D7&:&K(\gamma)=\sum_i(a_i^\dagger\gamma_i-\gamma_i^*a_i),
      \quad D(\gamma)=e^{K(\gamma)}.\\
L2&:&K(\gamma)^\dagger=-K(\gamma),\text{ hence }D(\gamma)\text{ is unitary.}\\
L3&:&D(\gamma)^\dagger a_nD(\gamma)=a_n+\gamma_n\text{ and similarly for }
      a_n^\dagger.\\
L4&:&[K(\alpha),K(\beta)]=\sigma(\alpha,\beta).\\
T1&:&\gamma\mapsto D(\gamma)\text{ is a unitary representation of }\mathsf D_I
      \text{ in }\mathcal A_I.\\
T2&:&\text{A faithful ordinary complex-scalar Hilbert representation of this
      Grassmann cocycle is not the theorem.}
\end{array}
\]

\subsection{The Concrete Target Algebra D1--D4}

Let \(I=\{1,\ldots,N\}\).  The finite CAR algebra
\(\operatorname{CAR}(I)\) is the complex unital \(*\)-algebra generated by odd
symbols \(a_i,a_i^\dagger\), \(i\in I\), subject to
\[
  \{a_i,a_j^\dagger\}=\delta_{ij}1,\qquad
  \{a_i,a_j\}=0,\qquad
  \{a_i^\dagger,a_j^\dagger\}=0.
\]
Here \(\{x,y\}=xy+yx\).  The adjoint satisfies
\((a_i)^\dagger=a_i^\dagger\) and \((a_i^\dagger)^\dagger=a_i\).

Let \(\Lambda_I\) be the Grassmann \(*\)-algebra generated by odd symbols
\(\gamma_i,\gamma_i^*\), \(i\in I\), with
\[
  \{\gamma_i,\gamma_j\}=0,\qquad
  \{\gamma_i^*,\gamma_j\}=0,\qquad
  \{\gamma_i^*,\gamma_j^*\}=0.
\]
The symbols \(\gamma_i\) and \(\gamma_i^*\) are independent generators related
by the chosen adjoint.  Following the source, adjunction reverses the order of
all odd quantities:
\[
  (x_1x_2\cdots x_k)^\dagger
  =
  x_k^\dagger\cdots x_2^\dagger x_1^\dagger
\]
when the \(x_r\) are fermionic operators or Grassmann variables.

Define the Grassmann-extended CAR algebra
\[
  \mathcal A_I=\operatorname{CAR}(I)\widehat\otimes\Lambda_I
\]
as the unital \(*\)-algebra generated by the two preceding sets, with the
additional cross anticommutation rules
\[
  \{\gamma_i,a_j\}=\{\gamma_i,a_j^\dagger\}
  =\{\gamma_i^*,a_j\}=\{\gamma_i^*,a_j^\dagger\}=0 .
\]
An element is even if it is a sum of monomials with even total fermionic
degree, and odd if it has odd total degree.  Even Grassmann elements commute
with all CAR generators.  This is the algebra in which Cahill--Glauber's
unitary displacement operators live.

\subsection{The Abstract Ray Algebra D5--D6}

A Grassmann displacement parameter is an \(I\)-tuple
\(\gamma=(\gamma_i)_{i\in I}\) of odd Grassmann elements, with adjoint tuple
\(\gamma^*=(\gamma_i^*)_{i\in I}\).  The additive law is componentwise.

For two parameters \(\alpha,\beta\), define the even Grassmann form
\[
  \sigma(\alpha,\beta)
  =
  \sum_{i\in I}(\beta_i^*\alpha_i-\alpha_i^*\beta_i)
\]
and its exponential
\[
  c(\alpha,\beta)=\exp\!\left(\frac12\sigma(\alpha,\beta)\right).
\]
The exponential is a finite polynomial because every Grassmann monomial is
nilpotent.

\paragraph{Lemma \(L1\).}
\(\sigma\) is even, central in \(\mathcal A_I\), skew-adjoint, and gives a
multiplicative two-cocycle \(c\):
\[
  c(\alpha,\beta)c(\alpha+\beta,\delta)
  =
  c(\beta,\delta)c(\alpha,\beta+\delta).
\]

\emph{Proof.}
Each summand \(\beta_i^*\alpha_i-\alpha_i^*\beta_i\) is a product of two odd
Grassmann elements, hence is even.  Even Grassmann elements commute with all
Grassmann generators and with all CAR generators, so \(\sigma\) is central.
The adjoint convention gives
\[
  (\beta_i^*\alpha_i)^\dagger=\alpha_i^*\beta_i,\qquad
  (\alpha_i^*\beta_i)^\dagger=\beta_i^*\alpha_i,
\]
so \(\sigma(\alpha,\beta)^\dagger=-\sigma(\alpha,\beta)\).  Bilinearity gives
\[
\begin{aligned}
&\sigma(\alpha,\beta)+\sigma(\alpha+\beta,\delta)\\
&\qquad =
  \sigma(\alpha,\beta)+\sigma(\alpha,\delta)+\sigma(\beta,\delta)\\
&\qquad =
  \sigma(\beta,\delta)+\sigma(\alpha,\beta+\delta).
\end{aligned}
\]
Because the exponents are central, exponentials multiply by adding exponents.
This proves the displayed two-cocycle identity.

Define the abstract Grassmann displacement ray algebra \(\mathsf D_I\) to be
the unital \(*\)-algebra over the even Grassmann scalars generated by symbols
\(\mathsf D(\gamma)\), one for each displacement parameter, with relations
\[
  \mathsf D(0)=1,\qquad
  \mathsf D(\alpha)\mathsf D(\beta)
  =
  \mathsf D(\alpha+\beta)c(\alpha,\beta),\qquad
  \mathsf D(\gamma)^\dagger=\mathsf D(-\gamma).
\]
Lemma \(L1\) is exactly the associativity check for the product relation.

\subsection{The Cahill--Glauber Representation D7--T1}

For a parameter \(\gamma\), define
\[
  K(\gamma)=\sum_{i\in I}(a_i^\dagger\gamma_i-\gamma_i^*a_i),
  \qquad
  D(\gamma)=e^{K(\gamma)}\in\mathcal A_I.
\]
Each term \(a_i^\dagger\gamma_i\) and \(\gamma_i^*a_i\) is even, so
\(K(\gamma)\) is even and the ordinary commutator is the relevant bracket.

\paragraph{Theorem \(T1\): representation to be proved.}
The assignment
\[
  \Pi:\mathsf D_I\longrightarrow\mathcal A_I^\times,\qquad
  \Pi(\mathsf D(\gamma))=D(\gamma)
  =
  \exp\!\left(\sum_i(a_i^\dagger\gamma_i-\gamma_i^*a_i)\right),
\]
is a unitary representation of the abstract Grassmann displacement ray algebra
in the Grassmann-extended finite CAR algebra.  Its product law is
\[
  D(\alpha)D(\beta)
  =
  D(\alpha+\beta)
  \exp\!\left(\frac12\sum_i(\beta_i^*\alpha_i-\alpha_i^*\beta_i)\right).
\]

\subsection{Conclusion}

The unitary representation of the abstract fermion displacement objects is:
\[
  \boxed{\mathsf D(\gamma)\mapsto
  \exp\!\left(\sum_i(a_i^\dagger\gamma_i-\gamma_i^*a_i)\right)}
\]
inside the Grassmann-extended finite CAR algebra.  It is unitary, it displaces
the concrete CAR generators by \(a_i\mapsto a_i+\gamma_i\) and
\(a_i^\dagger\mapsto a_i^\dagger+\gamma_i^*\), and it obeys the
Grassmann-valued ray product of Cahill--Glauber.  The derivations of unitarity,
the displacement identities, the product law, and the ordinary Hilbert-space
boundary are in Section~\ref{sec:grassmann-weyl-fermion-displacement-derivations}.
